% Budget Guard -- NeurIPS 2026 submission, final draft (2026-08-24).
% Reconstructed from the 2026-08-23 condensed 18-page revision with ONLY the agreed
% fixes applied:
%   1. Abstract + Conclusion: "on one split each" restored on the 18-of-19 claim.
%   2. Checklist Q1 "headlined" -> "headline"; Q7 "reported" -> "headline".
%   3. Appendix C.9: BAGEN no-licence sentence and library licence list restored.
%   4. Checklist restored to the official template format (Question/Answer/
%      Justification/Guidelines).
%   5. Bibliography restored to the verified full-author-list entries
%      (checked against primary sources 2026-08-19; see paper/CLAIMS_EVIDENCE.md).
% Requires in the same project: neurips_2026.sty, fig1_pipeline.pdf,
% fig2_coverage.pdf, fig3_transfer.pdf, fig4_economics.pdf (all in
% paper/submission/BudgetGuard_NeurIPS2026_source.zip).

\documentclass{article}

\PassOptionsToPackage{numbers,compress}{natbib}
% Submission form: no [final], no [preprint] -- anonymised, with line numbers.
\usepackage{neurips_2026}

\usepackage[utf8]{inputenc}
\usepackage[T1]{fontenc}
\usepackage[hidelinks]{hyperref}
\usepackage{url}
\usepackage{booktabs}
\usepackage{amsfonts}
\usepackage{amsmath}
\usepackage{amssymb}
\usepackage{nicefrac}
\usepackage{microtype}
\usepackage{graphicx}
\usepackage{xcolor}

\newcommand{\bg}{Budget Guard}

% NeurIPS checklist answer macros (from the official template).
\providecommand{\answerYes}[1][]{\textcolor{blue}{[Yes] #1}}
\providecommand{\answerNo}[1][]{\textcolor{orange}{[No] #1}}
\providecommand{\answerNA}[1][]{\textcolor{gray}{[N/A] #1}}

\title{\bg: Calibrated Resource Control for Long-Horizon LLM Agents}

\author{%
  Anonymous Author(s)\\
  Affiliation\\
  Address\\
  \texttt{email}
}

\begin{document}

\maketitle

\begin{abstract}
Long-horizon LLM agents can consume substantial tokens, time, and money before
recognising that a task is unlikely to succeed. Prior work on budget-aware agents shows
the promise of early stopping, but agents' own remaining-budget estimates are poorly
calibrated and cannot be trusted for control. We introduce \bg, a model-agnostic,
training-free wrapper that supplies two complementary controls at each trajectory
prefix: conformal prediction intervals on remaining resource need, and a separately
calibrated feasibility abort gate. The intervals answer how much may still be required;
the gate answers whether continued execution is still worthwhile; they share features
and data but are distinct controls with separate calibration targets and evaluation
metrics. For interval calibration, Claude Opus 4.7 rollouts from Sokoban, Search-R1,
and SWE-bench reach 0.90--0.96 empirical coverage offline (split-CQR) and 0.87--0.92
online (windowed-CQR) against a nominal 0.90, where agent self-reports reach only 0.51,
0.23, and 0.35; across 19 model--environment runs, on one split each, windowed-CQR
holds at least 0.85 coverage in 18 of 19, and in multi-resource Warehouse
in-distribution Bonferroni joint calibration reaches 0.95 across time, capacity, and
cost, though it falls to 0.71 under a model swap. For abort economics, in Warehouse
only and without reading the remaining-budget intervals, thresholding the feasibility
classifier at $\delta{=}0.10$ avoids 86.3\% of doomed-run spend with no false aborts
observed in replay, and cuts Cost-of-Pass from \$1.61M to \$0.87M in simulated dollars
per successful run at full pass retention; at the same $\delta$, token-budget
environments save far less (7.1\% and 3.1\% of doomed token spend on Sokoban and
SWE-bench), and combining intervals into the abort rule is more conservative at fixed
$\delta$ and does not increase doomed-run savings. All stopping results are
counterfactual replays of logged rollouts, not live deployment. These results show
that reliable budget awareness can be supplied as an external statistical control
layer, without fine-tuning the underlying agent.
\end{abstract}

\section{Introduction}
\label{sec:intro}

LLM agents increasingly complete long-horizon tasks that require many reasoning steps,
tool calls, and interactions with external environments \cite{kwa2025long}, and those
capabilities carry growing resource demands \cite{chen2026token}. An agent may consume
thousands of tokens, substantial wall-clock time, or real money before finishing a task
\cite{bai2026spend}, and it may spend much of that on redundant reasoning that does not
improve the outcome \cite{sui2025overthinking,wang2025reasoning,yue2025dontoverthink}.
When a trajectory eventually fails, most of that spending is wasted \cite{bai2026spend}.
Controlling resource consumption without giving up task performance is therefore a
practical obstacle to deploying agentic systems at scale
\cite{chen2024frugalgpt,zellinger2025cascades,davidov2026quit}.

One difficulty is that an agent does not reliably know how much budget it still needs
\cite{bai2026spend}. It can produce an estimate, but those estimates are uncalibrated,
so they cannot be trusted for resource-control decisions.

This motivates \bg, a model-agnostic framework for calibrated resource control in
long-horizon LLM agents. Rather than teaching the underlying LLM budget awareness
through fine-tuning, \bg{} operates as an external statistical layer around an existing
agent. At each trajectory prefix it uses observable information (resources consumed so
far, progress, recent usage patterns, and the agent's own estimate when one is
available) to support two decisions. No LLM is trained or fine-tuned; the only fitted
components are lightweight gradient-boosted trees over usage-accounting features.

Rather than asking the agent for a single point estimate of remaining budget, \bg{}
produces an interval describing how much additional resource may still be required.
Conformal prediction calibrates these intervals so that they achieve a desired
empirical coverage rate under stated assumptions (exchangeability offline; rollout-safe
online evaluation when labels arrive only at trajectory completion). Online
recalibration absorbs temporal dependence and deployment shift when the underlying
model or task type changes. Separately, a feasibility classifier, calibrated at the
rollout level, supports an early-stopping rule with a user-chosen false-abort tolerance
$\delta$. These two components share features and data but serve different roles: the
intervals answer ``how much might remain?''; the gate answers ``should we stop?'' We
evaluate interval calibration and abort economics independently, each against metrics
matched to its role.

Our contributions are fourfold. (1)~We introduce a conformal framework for calibrating
progressive remaining-budget estimates that requires no LLM fine-tuning
(\S\ref{sec:method}, \S\ref{sec:within}; Table~\ref{tab:a2}). (2)~We evaluate adaptive
and sequential calibration under environment and model shift, where frozen calibration
collapses from 0.92 to 0.10 coverage and online recalibration recovers much of it
(\S\ref{sec:shift}; Figure~\ref{fig:transfer}). (3)~We extend budget awareness to
jointly constrained resources and censored trajectories (\S\ref{sec:method},
\S\ref{sec:conditional}, \S\ref{sec:censoring}; Table~\ref{tab:multires}). (4)~We
evaluate the feasibility abort gate independently of the remaining-budget intervals,
reporting Warehouse abort performance together with token-environment results at the
same $\delta$ and a classifier+interval variant that is more conservative and saves
less at that operating point (\S\ref{sec:econ}; Figure~\ref{fig:econ},
Table~\ref{tab:abort}).

\bg{} treats budget awareness as both a measurement problem and a control problem; the
interval and abort signals are distinct and are calibrated and reported on separate
metrics throughout.

\section{Related work}
\label{sec:related}

BAGEN \cite{lin2026bagen} formalises the estimation of an agent's remaining budget as
budget awareness and evaluates whether an agent can estimate its remaining requirements
mid-execution. A point estimate says nothing about its own uncertainty, and an interval
is only useful if its coverage can be trusted; BAGEN takes intervals directly from the
agent, and they are uncalibrated. Conformal prediction, by contrast, constructs
intervals with specified coverage guarantees without distributional assumptions
\cite{angelopoulos2023gentle,lei2018distribution}. BAGEN also asks the agent to judge,
at different prefixes, whether successful completion is still feasible, and covers
internal resources such as tokens alongside external ones such as time, money, and
warehouse capacity across Sokoban, Search-R1, SWE-bench, and Warehouse.

The BAGEN results expose a gap between task capability and budget awareness: strong
agents solve hard tasks while remaining poorly calibrated about what finishing will
cost. BAGEN reports that early stopping on budget-awareness signals saves roughly
28--64\% of tokens on failed trajectories, at a cost of roughly 1.6--4.2 percentage
points of task success, with interval coverage reaching only 47\% even after supervised
fine-tuning and reinforcement learning \cite{lin2026bagen}. The agent carries some
signal about future requirements, but its raw estimates are not reliable enough for
high-confidence control.

Doomed From the Start \cite{ruan2026doomed} monitors a model's internal signals to
decide when a trajectory is doomed and terminates it early. In common with \bg{} and
unlike BAGEN, it applies a distribution-free calibrated abort gate with no fine-tuning,
meeting recall targets of 90--97\% on TextCraft and saving 47.1\% (Qwen-2.5-7B) and
37.2\% (Llama-3.2-3B) of inference compute at the 90\% target. Because it reads
internal states, it applies only to open-weight models.

A third line applies conformal prediction over an agent's answer. MiCP
\cite{zhou2026micp} allocates an error budget to each turn of a multi-turn reasoning
process, letting an agent stop once it has gathered enough evidence while preserving an
overall coverage guarantee. MiCP and \bg{} share conformal machinery but calibrate
different objects: MiCP calibrates coverage over the answer; \bg{} calibrates coverage
over the remaining resource budget. \bg{} is also an external wrapper rather than a
component of the reasoning loop, and additionally handles jointly constrained
resources, censored trajectories, and a feasibility abort gate that we evaluate
separately from the intervals (\S\ref{sec:econ}). We specify the architecture in
\S\ref{sec:method} (Figure~\ref{fig:pipeline}); Appendix~\ref{app:background} defines
the background concepts used to interpret our results.

\section{Method}
\label{sec:method}

\paragraph{Setup.} At turn $k$ of a rollout, \bg{} observes only the prefix: per-turn
resource usage so far and, when available, the agent's own estimate of its remaining
budget. It outputs an interval for $y$, the resources the run still needs to finish
successfully, at nominal miscoverage $\alpha{=}0.10$. Two properties of this label
shape the design. First, $y$ is observable only once the rollout finishes, so no online
method may consume a label mid-run. Second, $y$ is defined only while the run can still
finish: a doomed run censors its label rather than delaying it. The guard acts as a
verifier: it need not trust the agent's self-report, because it wraps that self-report
in a guarantee that holds regardless of its quality (Table~\ref{tab:a1}).

\begin{figure}[!htb]
  \centering
  \includegraphics[width=\linewidth]{fig1_pipeline.pdf}
  \caption{\bg{} wraps an unmodified agent in an external calibration layer. Stages 1--4
  share prefix features; stage 3 returns a calibrated remaining-budget interval
  (\S\S\ref{sec:within}--\ref{sec:conditional}), and a separately calibrated
  feasibility score drives continue-or-abort (\S\ref{sec:econ}). Warehouse abort
  economics use the feasibility gate alone and do not read the intervals. Labels become
  available only when a rollout finishes, so online recalibration never uses a
  still-running run.}
  \label{fig:pipeline}
\end{figure}

Figure~\ref{fig:pipeline} sketches the two paths (interval calibration and feasibility
abort), which share prefix features but are evaluated separately in what follows. The
rest of this section specifies offline and online interval calibration for deployment
shift, joint calibration when several resources bind at once, censoring-aware bounds
behind a feasibility gate, and the economic stopping rule built on the gate alone.

\paragraph{Prefix features.} We use four families of features, all computable at turn
$k$ with no access to the rest of the trajectory. Three come from the usage stream
alone: progress (the turn index $k$); cumulative consumption (spend so far, and the
burn rate $\mathrm{spend}/k$); and recent-usage statistics (the last, mean, standard
deviation, minimum and maximum per-turn spend, the mean of the last three turns, and
the difference between the last turn and the mean). When available, a fourth family
encodes the agent's self-estimate (interval endpoints, midpoint, and width). Features
that depend on the rollout's future (final totals, success, and feasibility) are
excluded from the feature matrix and asserted absent in code. The guard needs no model
weights, logits, or environment internals: every feature comes from usage accounting
that provider APIs already report.

\paragraph{Calibrated intervals under deployment shift.} The base predictor is a
gradient-boosted quantile regressor \cite{ke2017lightgbm}. We split calibration and
test by rollout, never by prefix: prefixes of one rollout share its final total, so
splitting them apart would leak each test rollout's answer into calibration. Frozen
split-CQR holds only under exchangeability, which breaks under a model swap or a new
task type. We therefore also recalibrate from the deployment stream's own finished
rollouts three ways: ACI, which raises or lowers the working coverage level from recent
misses \cite{gibbs2021aci}; DtACI, which removes ACI's hand-set step size
\cite{gibbs2024dtaci}; and windowed-CQR, a conformal quantile over a sliding window of
the 100 most recent scores, inspired by SPCI \cite{xu2023spci}; under
non-exchangeability the coverage gap is governed by window-to-stream drift
\cite{barber2023beyond}. All observe labels rollout-safely: a rollout's labels enter
the update state only after it finishes, and the state is frozen within a rollout.
Revealing them immediately would use labels that do not yet exist at deployment and
would inflate coverage, so we report only this deployment-faithful protocol. Two
reporting rules follow. An unbounded upper endpoint covers everything and informs
nothing, so we report the fraction of unbounded intervals beside every coverage number.
And because a single global quantile is set by the many easy late-turn prefixes, a good
average can hide under-coverage at the first turns, where an abort decision has the
most leverage; we therefore also calibrate a separate conformal quantile per turn-depth
bucket (Mondrian \cite{vovk2012conditional}): $k{=}1$, $k{=}2$, $k{\geq}3$.

\paragraph{Joint multi-resource calibration.} When a run is constrained by several
budgets at once (Warehouse: weeks, warehouse item-weeks, and dollars), the relevant
event is joint: the run survives only if all three hold, and covering each at level
$\alpha$ does not cover their conjunction, because the misses accumulate. Calibrating
each resource at $\alpha/3$ (Bonferroni) restores a joint guarantee of $1{-}\alpha$ by
the union bound, at the price of wider margins; this is our primary multi-resource
calibration scheme. A tighter alternative shares one threshold by ranking the
resources' errors jointly \cite{messoudi2021copula}, but that rank is chosen on the
same scores it then thresholds, so it carries no guarantee; we report it as an
observation and rest the claim on Bonferroni. Under a model swap the joint calibration
also moves: a rolling window of recent three-resource errors updates both the margins
and the shared rank, again only once a run has finished.

\paragraph{Censored rollouts and a selective guarantee.} When a run is doomed we never
learn $y$, only that it exceeded what was left. We calibrate one-sided upper bounds
five ways: treating censored scores as $+\infty$ (valid, but vacuous once the censored
fraction exceeds $\alpha$); substituting the known lower bound; dropping censored
scores entirely (the naive comparison); reweighting by each run's censoring probability
(IPCW \cite{gui2024survival,robins1992ipcw}); and charging a censored run only against
the resource that stopped it. A feasibility gate, a classifier on the same prefix
features, then issues intervals only on predicted-feasible prefixes, so coverage holds
among the prefixes it admits rather than across all prefixes. This selective guarantee
is what the construction supports; marginal coverage over all prefixes is reported as
an empirical observation.

\paragraph{Economic stopping.} The abort rule thresholds feasibility: stop at the first
turn whose feasibility score falls below $\tau$, the abort threshold. Because the rule
aborts a rollout exactly when that rollout's minimum prefix score crosses $\tau$, we
calibrate at the rollout level, setting $\tau$ to the $\delta$-quantile (with the
finite-sample correction) of the minimum scores of feasible calibration rollouts. This
bounds by $\delta$ (the false-abort tolerance) the probability of aborting a run that
would have finished: a per-run guarantee matched to the sequential decision, not a
per-prefix statement. Aborting a doomed run recovers its remaining spend, while
aborting a feasible run forfeits its pass; both effects enter Cost-of-Pass
\cite{erol2026costofpass}, total spend per passing run, so the rule is evaluated in
economic terms rather than as a coverage statistic. The rule is a calibrated binary
gate in the sense of conformal risk control \cite{angelopoulos2024risk}: it thresholds
the feasibility score and does not read the remaining-budget intervals. An
interval-driven variant that aborts when spend-so-far plus the calibrated lower bound
exceeds the budget is implemented in the released code; Warehouse economics below use
the feasibility-gate rule only, and token environments also evaluate the combined
variant (Table~\ref{tab:abort}).

\section{Experimental setup}
\label{sec:setup}

\begin{table}[t]
  \caption{The four evaluation environments; only Warehouse binds several resources at
  once, and only Warehouse censors labels. Token budgets are the BAGEN context budgets
  in tokens; Warehouse budgets are set per task. In Warehouse, half of the prefixes
  (640 of 1,280 per model) belong to runs that can no longer finish and drive the
  censored-calibration evaluation.}
  \label{tab:envs}
  \centering
  \small
  \begin{tabular}{lllll}
    \toprule
    Environment & Task & Budgeted resource(s) & Budget & \# Models \\
    \midrule
    Sokoban   & box-pushing puzzle          & tokens               & 2,500 tokens & 7 \\
    Search-R1 & multi-turn retrieval QA     & tokens               & 3,500 tokens & 7 \\
    SWE-bench & repository-level bug fixing & tokens               & 6,772--21,711 (per run) & 5 \\
    Warehouse & multi-week management sim.  & time, capacity, cost & per task, joint & 5 \\
    \bottomrule
  \end{tabular}
\end{table}

\paragraph{Data and environments.} All experiments use logged rollouts from BAGEN
\cite{lin2026bagen} at a pinned dataset revision; Table~\ref{tab:envs} summarises the
four environments. Sokoban, Search-R1, and SWE-bench are single-resource token tasks.
Sokoban and Search-R1 use uniform context budgets of 2,500 and 3,500 tokens; on
SWE-bench the official budget varies by run (6,772, 7,060, 8,265, 10,472, and 21,711
tokens across the five models), and we read each run's value from its estimation record
rather than applying a per-environment constant. The label is the number of tokens
still needed to complete the task under BAGEN's turn-excluding-history accounting.
Warehouse is a multi-week management simulation with three jointly binding budgets
(time in weeks, capacity in item-weeks, and cost in simulated USD), and per-prefix
ground truth for all three comes from the official estimation records. Half of its
prefixes (640 of 1,280 per model) belong to runs that can no longer finish, which makes
Warehouse our censoring test.

\paragraph{Protocol and baselines.} Rollouts are split 50/25/25 into train,
calibration, and test sets. We repeat every primary experiment with five random splits
(seeds 0--4), assigning whole rollouts to a split rather than individual prefixes, so
no rollout has its prefixes divided across sets; reported coverage and economics are
means over the five seeds. Online methods consume the test stream in its natural order
with labels revealed only at rollout boundaries. We evaluate three shift scenarios:
cross-environment (calibrate on Sokoban; deploy on Search-R1 and SWE-bench), model swap
(calibrate on Claude Opus 4.7; deploy on four other agents' SWE-bench streams), and an
external-resource model swap on Warehouse. Seven comparators run on the same splits and
seeds as the guard: the agent's own interval, a feature-free constant-quantile
interval, split conformal on absolute residuals, a conformalized-agent interval, a
guard-plus-self-estimate ablation, frozen split-CQR under shift, and hard caps for
stopping. Coverage is reported against the 0.90 nominal level, always alongside the
fraction of unbounded intervals. Appendix~\ref{app:protocol} gives the full protocol,
baseline definitions, metric formulas, and implementation details.

\paragraph{Models.} We evaluate seven agents (Claude Opus 4.7/4.6, Claude Sonnet 4.6,
GPT-5.2-instant, GPT-5.2-Codex, Gemini 3.1 Pro, and Qwen3-235B). Crossed with
environments this gives 19 token-budget model$\times$environment runs (7 Sokoban,
7 Search-R1, 5 SWE-bench) of 52--128 rollouts and 85--3,715 prefix samples each.

\section{Results}
\label{sec:results}

We evaluate prediction intervals using empirical coverage, normalised width, and the
fraction of unbounded intervals. Unbounded intervals raise coverage while carrying no
information, so we hatch them in the figures and report the unbounded fraction beside
every coverage number.

\subsection{Within-run interval calibration}
\label{sec:within}

\begin{figure}[!htb]
  \centering
  \includegraphics[width=\linewidth]{fig2_coverage.pdf}
  \caption{Conformal calibration closes almost all of the gap between agent
  self-reports and the 0.90 target, and windowed-CQR does so while maintaining bounded
  intervals throughout. Marginal coverage of the true remaining budget on Sokoban,
  Search-R1, and SWE-bench (Claude Opus 4.7 rollouts), $\alpha{=}0.10$, averaged over
  five rollout-level split seeds; error bars are the standard deviation across those
  seeds. Hatched bars emit unbounded intervals on the annotated fraction of the stream,
  which raises coverage at the cost of informativeness. Agent self-report is a single
  value read from the logged records rather than a seed mean. Coverage, normalised
  width, unbounded fraction, and test-set sizes for all methods are in
  Table~\ref{tab:a2}.}
  \label{fig:coverage}
\end{figure}

Figure~\ref{fig:coverage} shows within-run remaining-budget calibration on Sokoban,
Search-R1, and SWE-bench using Claude Opus 4.7 rollouts. Agent self-reported intervals
are badly calibrated, covering only 0.51, 0.23, and 0.35 of prefixes against a 0.90
target. Offline conformal calibration largely closes that gap, with split-CQR reaching
0.92, 0.90, and 0.96. Online windowed-CQR reaches 0.92, 0.92, and 0.87 with no
unbounded intervals in any of the three environments. ACI-CQR and DtACI sometimes reach
higher coverage, but do so partly by emitting unbounded intervals on up to 16\% of the
stream, which shows that coverage alone is not a sufficient criterion. Across all 19
logged model--environment runs, windowed-CQR reaches at least 0.85 coverage in 18 of 19
and at least 0.90 in 11 of 19, against 14 and 8 for DtACI, on one split per run
(Table~\ref{tab:b1}). It is also sharper than the feature-free baseline on two of three
environments: 22\% narrower than constant-quantile on Sokoban and 55\% narrower on
SWE-bench, but 23\% wider on Search-R1, where constant-quantile is the sharpest bounded
method (Table~\ref{tab:a2}). Adding the agent's own estimate to the guard's features
does not help: guard+self-est is within 0.01 of split-CQR on Sokoban (0.91 against
0.92), identical on SWE-bench (0.96), and worse on Search-R1 (0.85 against 0.90).

\subsection{Robustness under deployment shift}
\label{sec:shift}

\begin{figure}[!htb]
  \centering
  \includegraphics[width=\linewidth]{fig3_transfer.pdf}
  \caption{An environment change collapses frozen calibration and online recalibration
  recovers much of it; a model swap leaves every method below nominal coverage.
  Marginal coverage on the target stream. (a)~A guard calibrated on Sokoban and
  deployed unchanged on Search-R1 and SWE-bench: frozen split-CQR falls to 0.51 and
  0.10, and windowed-CQR recovers to 0.71 and 0.80. (b)~A guard calibrated on Claude
  Opus 4.7 SWE-bench rollouts and deployed on four other agents: here frozen and online
  calibration are comparable on average (0.74 against 0.73) and neither reaches the
  target, though windowed-CQR varies far less across seeds. Hatched bars emit unbounded
  intervals on the annotated fraction of the stream (60--88\% for the ACI-based methods
  in (a), and 22--61\% (ACI-CQR) and 4--7\% (DtACI) in (b)).}
  \label{fig:transfer}
\end{figure}

Figure~\ref{fig:transfer} evaluates \bg{} under deployment shift. Transferring a guard
calibrated on Sokoban to Search-R1 and SWE-bench drops frozen split-CQR coverage to
0.51 and 0.10. Rollout-safe windowed-CQR partially recovers it, to 0.71 and 0.80, with
no unbounded intervals; ACI and DtACI reach higher coverage but need unbounded
intervals on 60--88\% of predictions to get there. In the model-swap setting a guard
calibrated on Claude Opus 4.7 is transferred to Sonnet, Gemini, GPT, and Qwen SWE-bench
streams. Windowed-CQR reaches 0.66--0.89 there, and its seed-to-seed standard deviation
is an order of magnitude smaller than frozen calibration's (0.006--0.011 against
0.122--0.224), so its failures are at least predictable. Online adaptation improves
robustness under shift, but severe shift remains unsolved: no bounded-interval method
recovers the nominal level, and high coverage has to be read together with interval
informativeness.

\subsection{Conditional and multi-resource calibration}
\label{sec:conditional}

Marginal coverage can mask under-coverage in conditional slices. On Sokoban, split-CQR
reaches 0.92 marginal coverage but only 0.86 at the first turn, exactly where an abort
decision has the most leverage; Mondrian calibration by turn depth raises first-turn
coverage to 0.93. On Search-R1, Mondrian is neutral to slightly negative relative to
split-CQR; on SWE-bench the marginal-to-first-turn gap is already small and Mondrian is
neutral to slightly positive (Table~\ref{tab:b2}), yielding conditional validity rather
than a uniform improvement.

\begin{table}[!htb]
  \caption{Independent margins under-cover the joint event, Bonferroni fixes it in
  distribution, and nothing with a guarantee survives a model swap. Warehouse coverage
  against a nominal 0.90, mean over five rollout-level split seeds; above 0.90,
  coverage is bought with wider margins. Per-resource columns and the in-distribution
  joint column are Claude Opus 4.7; the final column is the joint event averaged over
  four swap targets. Bold marks the highest joint coverage in each column. ``--'' marks
  a scheme not run in that setting.}
  \label{tab:multires}
  \centering
  \small
  \begin{tabular}{lccccc}
    \toprule
    & \multicolumn{3}{c}{per-resource} & \multicolumn{2}{c}{joint} \\
    \cmidrule(lr){2-4} \cmidrule(lr){5-6}
    Calibration scheme & time & capacity & cost & in-dist. & model swap \\
    \midrule
    agent self-report & 0.87 & 0.67 & 0.71 & 0.42 & 0.34 \\
    independent margins ($\alpha$) & 0.99 & 0.89 & 0.92 & 0.80 & -- \\
    Bonferroni ($\alpha/3$) & 1.00 & 0.98 & 0.97 & \textbf{0.95} & 0.71 \\
    empirical copula & 0.99 & 0.94 & 0.96 & 0.89 & 0.62 \\
    adaptive copula (online) & -- & -- & -- & -- & \textbf{0.87} \\
    \bottomrule
  \end{tabular}
\end{table}

Table~\ref{tab:multires} reports the multi-resource case. In Warehouse, agent
self-reports reach 0.42 joint coverage across time, capacity, and cost. Calibrating
each resource independently at $\alpha$ lifts the joint event only to 0.80, reflecting
that independent margins do not compose to joint coverage. Bonferroni at $\alpha/3$
restores 0.95, and the empirical copula reaches 0.89 with narrower margins but no
guarantee. Under an external-resource model swap, however, frozen Bonferroni falls to
0.71, below its nominal level, and below the 0.87 reached by the un-guaranteed adaptive
copula, which recalibrates from finished rollouts. We report both the guaranteed and
empirical methods.

\subsection{Economic impact}
\label{sec:econ}

\begin{figure}[!htb]
  \centering
  \includegraphics[width=\linewidth]{fig4_economics.pdf}
  \caption{Warehouse economic impact from the feasibility gate alone, without
  remaining-budget intervals. At $\delta{=}0.10$ the classifier-only rule avoids 86.3\%
  of doomed-run spend with no observed false aborts, which no hard cap achieves at any
  setting. Simulated USD, mean over five seeds. (a)~Every hard cap that saves more than
  0.2\% of doomed spend also stops runs that would have succeeded; the gate occupies
  the favourable region of the savings--false-abort trade-off. (b)~Cost-of-Pass against
  the fraction of successful runs retained: the gate roughly halves cost per pass while
  retaining every observed success, whereas no cap improves on the unstopped baseline.
  Both panels are counterfactual replays of logged rollouts, not live deployments.}
  \label{fig:econ}
\end{figure}

Figure~\ref{fig:econ} evaluates early stopping in Warehouse using the classifier-only
feasibility gate. At $\delta{=}0.10$ the calibrated gate avoids 86.3\% of doomed-run
spend with 0\% observed false aborts. Hard spending caps achieve almost nothing at zero
false aborts (the tightest cap that preserves every successful run saves 0.2\% of
doomed spend), and dollar caps reach meaningful savings only at false-abort rates of
19--55\%; the only week cap that saves anything (54.5\%) aborts every feasible run. In
the Cost-of-Pass analysis the guard reduces estimated cost per successful trajectory
from \$1.61M to \$0.87M, a reduction of 46\%, while retaining all observed successful
runs. No hard cap improves Cost-of-Pass relative to the unstopped baseline by more than
a rounding error: discarding passes raises cost per pass, to \$3.19M at the tightest
cap.

\begin{table}[!htb]
  \caption{Abort performance at $\delta{=}0.10$ (mean over five seeds). Warehouse uses
  the classifier-only gate (native USD). Token abort is evaluated on Sokoban and
  SWE-bench only; both also evaluate a classifier+interval variant. Combining intervals
  reduces doomed savings at this operating point.}
  \label{tab:abort}
  \centering
  \small
  \begin{tabular}{llccc}
    \toprule
    Environment & Rule & False abort & Doomed spend saved & Uses intervals? \\
    \midrule
    Warehouse & feasibility only     & 0.0\% & 86.3\% & No \\
    Sokoban   & feasibility only     & 5.8\% & 7.1\%  & No \\
    Sokoban   & feasibility+interval & 5.8\% & 1.9\%  & Yes \\
    SWE-bench & feasibility only     & 0.0\% & 3.1\%  & No \\
    SWE-bench & feasibility+interval & 0.0\% & 0.0\%  & Yes \\
    \bottomrule
  \end{tabular}
\end{table}

Table~\ref{tab:abort} reports abort performance across environments and rules. At the
same $\delta{=}0.10$, Sokoban and SWE-bench save only 7.1\% and 3.1\% of doomed token
spend under the classifier-only rule, with false-abort rates of 5.8\% and 0\%.
Combining the interval trigger makes the policy more conservative at fixed $\delta$
(Sokoban savings fall to 1.9\%; SWE-bench to 0\%) and does not recover Warehouse-scale
savings. The 86.3\% figure is therefore a Warehouse result for the feasibility gate
alone, not a general \bg{} savings rate, and it does not use conformal remaining-budget
intervals. Two further cautions: the dollars are simulated Warehouse dollars, not
comparable to BAGEN's 28--64\% token savings; and the observed 0\% false-abort rate
should be read as consistent with the $\delta{=}0.10$ bound on a small test set, not as
proof that the true rate is zero.

\subsection{Censored calibration}
\label{sec:censoring}

Roughly half of Warehouse prefixes are censored: the run can no longer finish, so its
true completion requirement is never observed. Dropping those observations yields only
0.50 coverage on the censored prefixes where cost is the binding resource, so failed
trajectories cannot be treated as missing data. Treating censored scores as $+\infty$
is valid but vacuous at this censoring rate, and the informative treatments trade
validity for sharpness in ways the appendix makes explicit. Feasibility-gated selective
intervals offer the other lever: abstaining on predicted-doomed prefixes holds coverage
at 0.92--0.95 for capacity and cost, and at 1.00 for time, on the prefixes it keeps, at
the cost of keeping between 51\% and 89\% of them (Table~\ref{tab:b3}). The gate is
swept over its own prefix-level tolerance $\delta_{\mathrm{gate}}$, which is distinct
from the rollout-level $\delta$ of \S\ref{sec:econ}. One caveat governs how these
numbers should be read: the recovered labels used to score censored prefixes are exact
only for time, and inherit a week-proportional interpolation with 12--27\% median
relative error for capacity and cost, so those two are coverage against an approximate
label rather than ground truth.

\section{Discussion and limitations}
\label{sec:limits}

Our joint results rest entirely on Warehouse. Bonferroni calibration reaches 0.95 joint
coverage there, but a single multi-resource environment cannot support a claim that the
approach generalises, and \S\ref{sec:conditional} shows that guarantee is not robust to
a model swap. The 86.3\% doomed-spend figure likewise does not transfer to token
environments at $\delta{=}0.10$ (Table~\ref{tab:abort}).

Our breadth figures are single-split counts: on seed 0, windowed-CQR reaches at least
0.85 coverage in 18 of 19 runs, not all 19. Search-R1 leaves about 24 test prefixes per
seed, which is why its standard deviations reach 0.10--0.17. Coverage alone does not
separate a useful interval from a trivially wide one; we therefore report width and
unbounded fraction throughout. Censored capacity/cost results are measured against an
interpolated label. Finally, savings and Cost-of-Pass are counterfactual replays: they
assume adding the guard would not change agent behaviour and that the sample is
representative.

\paragraph{Broader impact.} The intended effect is to reduce wasted compute, energy,
and money on agent runs that cannot succeed; the Warehouse feasibility gate at
$\delta{=}0.10$ illustrates that potential in one environment (86.3\% of doomed-run
spend avoided in replay, without using remaining-budget intervals). The corresponding
risk is that a stopping threshold set too aggressively terminates runs that would have
finished; $\delta$ calibration exists to bound that risk at a user-chosen level, and
the counterfactual caveat above means a deployer should not treat the reported savings
as a guaranteed live return when choosing $\delta$, nor assume token environments will
see Warehouse-scale savings.

\section{Conclusion}
\label{sec:conclusion}

\bg{} shows that reliable budget awareness can be supplied as an external control layer
for long-horizon LLM agents, with no fine-tuning of the underlying model. Conformal
calibration of progressive remaining-budget estimates yields 0.90--0.96 empirical
coverage offline (split-CQR) and 0.87--0.92 online (windowed-CQR) against a 0.90 target
on Claude Opus 4.7 rollouts, and holds at least 0.85 coverage in 18 of 19
model$\times$environment runs, on one split each. Joint multi-resource calibration
works in distribution but does not survive a model swap. Separately, a calibrated
feasibility gate (not the remaining-budget intervals) avoids 86.3\% of doomed-run spend
in Warehouse replay at $\delta{=}0.10$, while token-budget environments yield much
smaller savings at the same operating point. Budget awareness can therefore serve not
only as a measurement but as a control signal during execution.

\begin{thebibliography}{27}

\bibitem{kwa2025long}
Thomas Kwa et~al.
\newblock Measuring {AI} ability to complete long software tasks.
\newblock In {\em Advances in Neural Information Processing Systems 38 (NeurIPS)},
  2025. arXiv:2503.14499.

\bibitem{chen2026token}
Yuxi Chen, Junming Chen, Chenyu He, Yiwei Li, Yicheng Ji, Yifan Wu, Dingyu Yang,
  Lansong Diao, Lidan Shou, Hongliang Zhang, Huan Li, and Gang Chen.
\newblock Token economics for {LLM} agents: A dual-view study from computing and
  economics.
\newblock {\em arXiv preprint arXiv:2605.09104}, 2026.

\bibitem{bai2026spend}
Longju Bai, Zhemin Huang, Xingyao Wang, Jiao Sun, Rada Mihalcea, Erik Brynjolfsson,
  Alex Pentland, and Jiaxin Pei.
\newblock How do {AI} agents spend your money? {A}nalyzing and predicting token
  consumption in agentic coding tasks.
\newblock {\em arXiv preprint arXiv:2604.22750}, 2026.

\bibitem{sui2025overthinking}
Yang Sui et~al.
\newblock Stop overthinking: A survey on efficient reasoning for large language
  models.
\newblock {\em Transactions on Machine Learning Research}, 2025. arXiv:2503.16419.

\bibitem{wang2025reasoning}
Rui Wang, Hongru Wang, Boyang Xue, Jianhui Pang, Shudong Liu, Yi Chen, Jiahao Qiu,
  Derek~Fai Wong, Heng Ji, and Kam-Fai Wong.
\newblock Harnessing the reasoning economy: A survey of efficient reasoning for large
  language models.
\newblock {\em arXiv preprint arXiv:2503.24377}, 2025.

\bibitem{yue2025dontoverthink}
Linan Yue, Yichao Du, Yizhi Wang, Weibo Gao, Fangzhou Yao, Li Wang, Ye Liu, Ziyu Xu,
  Qi Liu, Shimin Di, and Min-Ling Zhang.
\newblock Don't overthink it: A survey of efficient {R1}-style large reasoning models.
\newblock {\em arXiv preprint arXiv:2508.02120}, 2025.

\bibitem{chen2024frugalgpt}
Lingjiao Chen, Matei Zaharia, and James Zou.
\newblock {FrugalGPT}: How to use large language models while reducing cost and
  improving performance.
\newblock {\em Transactions on Machine Learning Research}, 2024. arXiv:2305.05176.

\bibitem{zellinger2025cascades}
Michael~J. Zellinger, Rex Liu, and Matt Thomson.
\newblock Cost-saving {LLM} cascades with early abstention.
\newblock {\em arXiv preprint arXiv:2502.09054}, 2025.

\bibitem{davidov2026quit}
Hen Davidov, Nachshon Cohen, Oren Kalinsky, Yaron Fairstein, Guy Kushilevitz, Ram
  Yazdi, and Patrick Rebeschini.
\newblock Knowing when to quit: A principled framework for dynamic abstention in {LLM}
  reasoning.
\newblock In {\em Proceedings of the 43rd International Conference on Machine Learning
  (ICML)}, 2026. arXiv:2604.18419.

\bibitem{romano2019cqr}
Yaniv Romano, Evan Patterson, and Emmanuel~J. Cand{\`e}s.
\newblock Conformalized quantile regression.
\newblock In {\em Advances in Neural Information Processing Systems 32 (NeurIPS)},
  pages 3538--3548, 2019. arXiv:1905.03222.

\bibitem{gibbs2021aci}
Isaac Gibbs and Emmanuel~J. Cand{\`e}s.
\newblock Adaptive conformal inference under distribution shift.
\newblock In {\em Advances in Neural Information Processing Systems 34 (NeurIPS)},
  2021. arXiv:2106.00170.

\bibitem{gibbs2024dtaci}
Isaac Gibbs and Emmanuel~J. Cand{\`e}s.
\newblock Conformal inference for online prediction with arbitrary distribution
  shifts.
\newblock {\em Journal of Machine Learning Research}, 25(162):1--36, 2024.
  arXiv:2208.08401.

\bibitem{lin2026bagen}
Yuxiang Lin, Zihan Wang, Mengyang Liu, Yuxuan Shan, Longju Bai, Junyao Zhang, Xing
  Jin, Boshan Chen, Jinyan Su, Xingyao Wang, Jiaxin Pei, and Manling Li.
\newblock {BAGEN}: Are {LLM} agents budget-aware?
\newblock {\em arXiv preprint arXiv:2606.00198}, 2026.

\bibitem{angelopoulos2023gentle}
Anastasios~N. Angelopoulos and Stephen Bates.
\newblock Conformal prediction: A gentle introduction.
\newblock {\em Foundations and Trends in Machine Learning}, 16(4):494--591, 2023.
  arXiv:2107.07511.

\bibitem{lei2018distribution}
Jing Lei, Max G'Sell, Alessandro Rinaldo, Ryan~J. Tibshirani, and Larry Wasserman.
\newblock Distribution-free predictive inference for regression.
\newblock {\em Journal of the American Statistical Association}, 113(523):1094--1111,
  2018. arXiv:1604.04173.

\bibitem{ruan2026doomed}
Kai Ruan, Zihe Huang, Ziqi Zhou, Qianshan Wei, Xuan Wang, and Hao Sun.
\newblock Doomed from the start: Early abort of {LLM} agent episodes via a
  recall-controlled probe cascade.
\newblock {\em arXiv preprint arXiv:2607.06503}, 2026.

\bibitem{zhou2026micp}
Xiaofan Zhou, Huy Nguyen, Bo Yu, Chenxi Liu, and Lu Cheng.
\newblock Adaptive stopping for multi-turn {LLM} reasoning.
\newblock {\em arXiv preprint arXiv:2604.01413}, 2026.

\bibitem{xu2023spci}
Chen Xu and Yao Xie.
\newblock Sequential predictive conformal inference for time series.
\newblock In {\em Proceedings of the 40th International Conference on Machine Learning
  (ICML)}, volume 202 of {\em PMLR}, pages 38707--38727, 2023. arXiv:2212.03463.

\bibitem{barber2023beyond}
Rina~Foygel Barber, Emmanuel~J. Cand{\`e}s, Aaditya Ramdas, and Ryan~J. Tibshirani.
\newblock Conformal prediction beyond exchangeability.
\newblock {\em The Annals of Statistics}, 51(2):816--845, 2023. arXiv:2202.13415.

\bibitem{vovk2012conditional}
Vladimir Vovk.
\newblock Conditional validity of inductive conformal predictors.
\newblock In {\em Proceedings of the Asian Conference on Machine Learning (ACML)},
  volume 25 of {\em PMLR}, pages 475--490, 2012. arXiv:1209.2673.

\bibitem{messoudi2021copula}
Soundouss Messoudi, S{\'e}bastien Destercke, and Sylvain Rousseau.
\newblock Copula-based conformal prediction for multi-target regression.
\newblock {\em Pattern Recognition}, 120:108101, 2021. arXiv:2101.12002.

\bibitem{gui2024survival}
Yu Gui, Rohan Hore, Zhimei Ren, and Rina~Foygel Barber.
\newblock Conformalized survival analysis with adaptive cut-offs.
\newblock {\em Biometrika}, 111(2):459--477, 2024. arXiv:2211.01227.

\bibitem{robins1992ipcw}
James~M. Robins and Andrea Rotnitzky.
\newblock Recovery of information and adjustment for dependent censoring using
  surrogate markers.
\newblock In Nicholas~P. Jewell, Klaus Dietz, and Vernon~T. Farewell, editors, {\em
  AIDS Epidemiology: Methodological Issues}, pages 297--331. Birkh{\"a}user Boston,
  Boston, MA, 1992.

\bibitem{kaplan1958nonparametric}
E.~L. Kaplan and Paul Meier.
\newblock Nonparametric estimation from incomplete observations.
\newblock {\em Journal of the American Statistical Association}, 53(282):457--481,
  1958.

\bibitem{erol2026costofpass}
Mehmet~Hamza Erol, Batu El, Mirac Suzgun, Mert Yuksekgonul, and James Zou.
\newblock Cost-of-pass: An economic framework for evaluating language models.
\newblock In {\em International Conference on Learning Representations (ICLR)}, 2026.
  arXiv:2504.13359.

\bibitem{angelopoulos2024risk}
Anastasios~N. Angelopoulos, Stephen Bates, Adam Fisch, Lihua Lei, and Tal Schuster.
\newblock Conformal risk control.
\newblock In {\em International Conference on Learning Representations (ICLR)}, 2024.
  arXiv:2208.02814.

\bibitem{ke2017lightgbm}
Guolin Ke, Qi Meng, Thomas Finley, Taifeng Wang, Wei Chen, Weidong Ma, Qiwei Ye, and
  Tie-Yan Liu.
\newblock {LightGBM}: A highly efficient gradient boosting decision tree.
\newblock In {\em Advances in Neural Information Processing Systems 30 (NeurIPS)},
  pages 3146--3154, 2017.

\end{thebibliography}

\appendix

\section{Background and guarantees}
\label{app:background}

This appendix defines the concepts needed to interpret the results.

\begin{itemize}
  \item \textbf{Coverage.} The fraction of cases in which the true remaining budget
  falls inside the predicted interval. Coverage of 0.90 means the true value is
  contained in about 90\% of cases.
  \item \textbf{Calibration set.} A held-out set of completed trajectories used to set
  interval widths and stopping thresholds.
  \item \textbf{Rollout (run).} One complete agent attempt at a task.
  \item \textbf{Turn ($k$).} One agent step within a rollout.
  \item \textbf{Prefix.} The partial trajectory containing all turns completed through
  turn $k$.
  \item \textbf{$\hat q$.} The conformal threshold: a learned error margin added to the
  model's predictions to expand them into an interval. Throughout, $\hat q$ is the
  conformal threshold and $\hat f_{\mathrm{lo}}$, $\hat f_{\mathrm{hi}}$ are the
  quantile regressor's own predictions.
  \item \textbf{$\alpha$.} The allowed miscoverage rate; the target coverage is
  $1{-}\alpha$. We use $\alpha{=}0.10$ throughout. A smaller $\alpha$ raises the target
  and generally widens intervals.
  \item \textbf{Width.} The size of the predicted range. Narrower is more informative
  at equal coverage: $[100, 115]$ is more useful than $[100, 300]$. An interval
  extending to infinity attains coverage while saying nothing.
  \item \textbf{$\delta$.} The false-abort tolerance for the stopping rule. A larger
  $\delta$ permits greater savings on doomed runs while accepting more risk of aborting
  successful ones.
  \item \textbf{$\tau$.} The abort threshold. The guard aborts at the first turn whose
  feasibility score falls below $\tau$. It is not chosen by hand: it is read off the
  calibration rollouts as the value holding the false-abort rate at or below $\delta$,
  so a smaller $\delta$ produces a more cautious $\tau$.
  \item \textbf{Offline.} Calibrates once on completed runs and stays fixed during
  deployment.
  \item \textbf{Online.} Updates calibration after each completed run, using its newly
  observed remaining-budget labels.
\end{itemize}

\begin{table}[t]
  \caption{What each component guarantees, and what it does not. Statements we do not
  have are listed as such, so a reader can tell a theorem from an observation.}
  \label{tab:a1}
  \centering
  \footnotesize
  \begin{tabular}{@{}p{0.16\linewidth}p{0.36\linewidth}p{0.40\linewidth}@{}}
    \toprule
    Method & Role in \bg{} & Guarantee and scope \\
    \midrule
    split-CQR & CQR intervals conformalized on a by-rollout calibration split; our
    offline default \cite{romano2019cqr}. & Marginal coverage $\geq 1{-}\alpha$ over a
    random (rollout, prefix) pair, under exchangeable rollouts. Pooling a rollout's
    prefixes is an empirical evaluation, not the guaranteed object
    \cite{lei2018distribution}. \\
    \addlinespace
    Mondrian-CQR & split-CQR calibrated separately per turn-depth bucket
    \cite{vovk2012conditional}. & Coverage $\geq 1{-}\alpha$ within each bucket, under
    exchangeability within the bucket; thin buckets fall back to the global quantile. \\
    \addlinespace
    ACI / DtACI & Adaptive coverage levels recalibrated from the stream's finished
    rollouts \cite{gibbs2021aci,gibbs2024dtaci}. & Long-run average coverage approaches
    $1{-}\alpha$. May emit unbounded intervals; we clamp the working level at 0.999
    rather than emit empty sets, a conservative deviation. \\
    \addlinespace
    windowed-CQR & Conformal quantile over a sliding window of the 100 most recent
    scores; an SPCI-inspired simplification \cite{xu2023spci}. & No finite-sample
    guarantee; coverage gap bounded by window-to-stream drift \cite{barber2023beyond}.
    Evaluated empirically. \\
    \addlinespace
    Bonferroni joint & Calibrates each resource at $\alpha/3$ so all hold at once; our
    primary multi-resource calibration scheme. & Joint coverage $\geq 1{-}\alpha$ by
    the union bound; no data-driven selection; conservative, so realised coverage
    exceeds the target in distribution. \\
    \addlinespace
    empirical copula & Shared copula rank across resources \cite{messoudi2021copula};
    an adaptive variant updates the rank per finished rollout. & Empirical only: the
    rank is selected on the same scores it thresholds. Reported alongside Bonferroni,
    not as the guaranteed claim. \\
    \addlinespace
    censoring-aware upper bounds & One-sided bounds for doomed runs
    \cite{gui2024survival,robins1992ipcw,kaplan1958nonparametric}. & One-sided upper
    coverage on the censored subset. $+\infty$ always valid but vacuous once the
    censored fraction exceeds $\alpha$; the others empirical or assumption-dependent. \\
    \addlinespace
    feasibility-gated intervals & A classifier on the same prefix features issues
    intervals only on predicted-feasible prefixes. & Coverage $\geq 1{-}\alpha$
    conditional on passing the gate: selective, not marginal. Marginal coverage
    reported empirically. \\
    \addlinespace
    stopping rule ($\tau$) & Abort at the first turn whose feasibility score drops
    below $\tau$ \cite{angelopoulos2024risk}; scored by Cost-of-Pass
    \cite{erol2026costofpass}. & False-abort probability $\leq \delta$ per run, from
    rollout-level calibration on feasible runs. \\
    \bottomrule
  \end{tabular}
\end{table}

Table~\ref{tab:a1} states, component by component, what is guaranteed and under what
assumption. Four rows carry no finite-sample guarantee at all (windowed-CQR, the
empirical copula, the censoring-aware bounds other than $+\infty$, and the marginal
rather than selective reading of the gated intervals), and every claim resting on those
is reported in this paper as an empirical observation.

\section{Per-run breadth results}
\label{app:breadth}

Table~\ref{tab:b1} lists every model$\times$environment run behind the breadth claims
in \S\ref{sec:within}; the ``18 of 19'' and ``11 of 19'' counts can be read straight
off the windowed-CQR column. Table~\ref{tab:b2} gives conditional coverage by turn
depth, Table~\ref{tab:b3} the feasibility-gate sweep, and Table~\ref{tab:b4} the
censored upper-bound calibrations. Every number is read from the generated result
tables, which one command regenerates and two checkers verify.

\begin{table}[t]
  \caption{Breadth across all 19 token-budget runs: windowed-CQR clears 0.85 on 18 of
  them and 0.90 on 11, while the agents' own intervals clear neither anywhere. Marginal
  coverage against nominal 0.90 on a single rollout-level split (seed 0). Four decimals are shown on the three guard
  columns so the two counts above are directly checkable; at three decimals
  SWE-bench/Gemini DtACI rounds to 0.850 and appears to clear 0.85 when it does not.
  ``--'' marks a run whose records carry no agent interval. $^{\dagger}$Warehouse is
  listed for dataset composition only; its 0.424 is joint coverage across three
  resources (Table~\ref{tab:multires}), not a token-budget marginal.}
  \label{tab:b1}
  \centering
  \footnotesize
  \begin{tabular}{llrrcccc}
    \toprule
    Environment & Model & Rollouts & Prefixes & Agent & split-CQR & windowed-CQR & DtACI \\
    \midrule
    Sokoban & Claude Opus 4.6 & 116 & 410 & 0.446 & 0.8077 & 0.8750 & 0.8846 \\
     & Claude Opus 4.7 & 123 & 355 & 0.513 & 0.8667 & 0.9000 & 0.9000 \\
     & Claude Sonnet 4.6 & 126 & 400 & 0.524 & 0.9623 & 0.9151 & 0.9528 \\
     & Gemini 3.1 Pro & 126 & 389 & 0.125 & 0.9425 & 0.9310 & 0.9310 \\
     & GPT-5.2-instant & 128 & 661 & 0.481 & 0.8176 & 0.8986 & 0.8716 \\
     & GPT-5.2-Codex & 128 & 592 & -- & 0.9933 & 0.9600 & 0.9467 \\
     & Qwen3-235B & 126 & 1{,}038 & 0.090 & 0.9268 & 0.9024 & 0.8862 \\
    \midrule
    Search-R1 & Claude Opus 4.7 & 66 & 87 & 0.233 & 1.0000 & 1.0000 & 1.0000 \\
     & Claude Sonnet 4.6 & 64 & 102 & 0.386 & 1.0000 & 0.9200 & 1.0000 \\
     & Claude Opus 4.6 & 52 & 90 & 0.478 & 0.9130 & 0.9130 & 0.9130 \\
     & GPT-5.2-instant & 125 & 324 & 0.198 & 0.8256 & 0.8605 & 0.8488 \\
     & GPT-5.2-Codex & 126 & 262 & 0.202 & 0.9649 & 0.9298 & 0.8772 \\
     & Gemini 3.1 Pro & 83 & 85 & 0.202 & 0.8182 & 0.8636 & 0.8182 \\
     & Qwen3-235B & 124 & 1{,}079 & 0.025 & 0.9259 & 0.9259 & 0.8765 \\
    \midrule
    SWE-bench & Claude Opus 4.7 & 64 & 712 & 0.354 & 0.9926 & 0.9185 & 0.9037 \\
     & Claude Sonnet 4.6 & 64 & 1{,}258 & 0.251 & 0.9101 & 0.8705 & 0.8957 \\
     & Gemini 3.1 Pro & 64 & 2{,}327 & 0.254 & 0.8497 & 0.8899 & 0.8497 \\
     & GPT-5.2-instant & 64 & 1{,}313 & 0.385 & 0.8840 & 0.8589 & 0.8370 \\
     & Qwen3-235B & 60 & 3{,}715 & 0.034 & 1.0000 & 0.7727 & 0.8271 \\
    \midrule
    Warehouse & Opus 4.7 $+$ 4 swap targets & 64 ea. & 1{,}280 ea. & 0.424$^{\dagger}$ & -- & -- & -- \\
    \bottomrule
  \end{tabular}
\end{table}

\begin{table}[t]
  \caption{Marginal coverage hides a first-turn deficit on Sokoban, and Mondrian
  calibration is what removes it. Coverage by turn-depth bucket, nominal 0.90, mean over five seeds. Search-R1
  rollouts are short enough that no $k{\geq}3$ bucket exists.}
  \label{tab:b2}
  \centering
  \small
  \begin{tabular}{llcccc}
    \toprule
    Environment & Method & Marginal & $k{=}1$ & $k{=}2$ & $k{\geq}3$ \\
    \midrule
    Sokoban & split-CQR & 0.92 & 0.86 & 0.93 & 0.96 \\
     & Mondrian-CQR & 0.94 & 0.93 & 0.96 & 0.94 \\
    \midrule
    Search-R1 & split-CQR & 0.90 & 0.91 & 0.89 & -- \\
     & Mondrian-CQR & 0.89 & 0.89 & 0.89 & -- \\
    \midrule
    SWE-bench & split-CQR & 0.96 & 0.99 & 0.99 & 0.95 \\
     & Mondrian-CQR & 0.97 & 0.99 & 0.99 & 0.96 \\
    \bottomrule
  \end{tabular}
\end{table}

\begin{table}[t]
  \caption{The feasibility gate trades retained volume for coverage on what it keeps:
  abstaining on predicted-doomed prefixes holds 0.92--0.95 coverage on capacity,
  where marginal calibration reaches only 0.89. Warehouse, mean over
  five seeds. ``Kept'' is the fraction of prefixes the gate admits, ``caught'' the
  fraction of doomed prefixes it correctly abstains on, ``false abort'' the fraction of
  feasible prefixes it wrongly rejects, and ``coverage'' the interval coverage among
  admitted feasible prefixes. Time binds in only 30\% of censored runs and the gate
  catches almost all of them, so its row is flat. $\delta$ here is the gate's
  prefix-level tolerance, not the rollout-level false-abort tolerance of
  \S\ref{sec:econ}.}
  \label{tab:b3}
  \centering
  \small
  \begin{tabular}{lccccc}
    \toprule
    Resource & $\delta$ & Kept ($\uparrow$) & Caught ($\uparrow$) & False abort ($\downarrow$) & Coverage ($\uparrow$) \\
    \midrule
    time (weeks) & 0.02 & 0.88 & 0.93 & 0.00 & 1.00 \\
     & 0.05 & 0.88 & 0.93 & 0.00 & 1.00 \\
     & 0.10 & 0.88 & 0.93 & 0.00 & 1.00 \\
     & 0.20 & 0.88 & 0.93 & 0.00 & 1.00 \\
     & 0.30 & 0.88 & 0.93 & 0.00 & 1.00 \\
    \midrule
    capacity (item-weeks) & 0.02 & 0.89 & 0.31 & 0.04 & 0.95 \\
     & 0.05 & 0.78 & 0.64 & 0.05 & 0.95 \\
     & 0.10 & 0.70 & 0.74 & 0.07 & 0.95 \\
     & 0.20 & 0.57 & 0.82 & 0.19 & 0.95 \\
     & 0.30 & 0.51 & 0.85 & 0.23 & 0.94 \\
    \midrule
    cost (USD) & 0.02 & 0.86 & 0.60 & 0.05 & 0.92 \\
     & 0.05 & 0.81 & 0.69 & 0.10 & 0.92 \\
     & 0.10 & 0.80 & 0.72 & 0.10 & 0.92 \\
     & 0.20 & 0.74 & 0.75 & 0.21 & 0.93 \\
     & 0.30 & 0.74 & 0.76 & 0.21 & 0.93 \\
    \bottomrule
  \end{tabular}
\end{table}

\begin{table}[t]
  \caption{Dropping censored runs under-covers exactly the expensive failures, and
  every informative repair trades validity for sharpness. One-sided upper-bound
  coverage on Warehouse at nominal 0.90, split by whether the resource was the binding
  constraint. $+\infty$ is always valid and always vacuous; ``drop'' is the naive
  comparison. Kaplan--Meier is a diagnostic, not a calibration option. Coverage for capacity and cost is measured
  against an interpolated label (12--27\% median relative error), not ground truth.}
  \label{tab:b4}
  \centering
  \small
  \begin{tabular}{llcccc}
    \toprule
    Resource & Treatment & Feasible & \multicolumn{2}{c}{Censored} & Binding \\
    \cmidrule(lr){4-5}
     & & ($\uparrow$0.90) & binding & non-binding & fraction \\
    \midrule
    time (weeks) & drop (naive) & 1.00 & 1.00 & 1.00 & 0.30 \\
     & $+\infty$ & 1.00 & 1.00 & 1.00 & 0.30 \\
     & substituted bound & 1.00 & 1.00 & 1.00 & 0.30 \\
     & IPCW & 0.96 & 0.96 & 0.96 & 0.30 \\
     & per-resource bound & 0.90 & 0.90 & 0.90 & 0.30 \\
     & Kaplan--Meier (diagnostic) & 1.00 & 1.00 & 1.00 & 0.30 \\
    \midrule
    capacity (item-weeks) & drop (naive) & 0.94 & 0.77 & 0.98 & 0.34 \\
     & $+\infty$ & 1.00 & 1.00 & 1.00 & 0.34 \\
     & substituted bound & 0.99 & 0.98 & 1.00 & 0.34 \\
     & IPCW & 0.93 & 0.75 & 0.98 & 0.34 \\
     & per-resource bound & 0.91 & 0.74 & 0.98 & 0.34 \\
     & Kaplan--Meier (diagnostic) & 0.94 & 0.78 & 0.99 & 0.34 \\
    \midrule
    cost (USD) & drop (naive) & 0.88 & 0.50 & 0.97 & 0.42 \\
     & $+\infty$ & 1.00 & 1.00 & 1.00 & 0.42 \\
     & substituted bound & 1.00 & 0.99 & 1.00 & 0.42 \\
     & IPCW & 0.88 & 0.50 & 0.96 & 0.42 \\
     & per-resource bound & 0.86 & 0.46 & 0.94 & 0.42 \\
     & Kaplan--Meier (diagnostic) & 0.92 & 0.58 & 0.98 & 0.42 \\
    \bottomrule
  \end{tabular}
\end{table}

\section{Protocol, baselines, metrics, and implementation}
\label{app:protocol}

This appendix records what \S\ref{sec:setup} summarises: how the data is split, what
each baseline is, how each reported number is computed, and the settings under which
all of it runs. Every value below is fixed in the released code, and the parenthetical
names are the modules that implement it.

\subsection{Splits and seeds}
\label{app:splits}

Every split is by rollout, never by prefix (\texttt{split.group\_split}). For a given
seed the unique rollout indices are permuted with
\texttt{numpy.random.default\_rng(seed)} and allocated 50/25/25 to train, calibration
and test; every prefix follows its rollout. Prefixes of one rollout share that
rollout's final total, so splitting them apart would place each test rollout's answer
into calibration. A split that is requested but comes out empty raises an error rather
than degrading quietly: with an empty calibration set the conformal quantile of zero
scores is $+\infty$, every interval becomes $[0, \infty)$, and coverage reads a
meaningless 1.00 with no other symptom. Reported numbers are means over five seeds,
0 through 4.

\subsection{Rollout-safe online evaluation}
\label{app:online}

Online methods consume the test rows in their natural trajectory order, rollout by
rollout and turn by turn. A prefix's label is the resource the run still needs to
finish, which cannot be observed until the rollout ends, so a rollout's labels enter
the update state only at its boundary and the state is frozen within a rollout
(\texttt{conformal.aci\_cqr}, \texttt{windowed\_cqr}, \texttt{dtaci}, each taking a
\texttt{groups} argument). The code rejects test streams whose rollout ids are not
contiguous blocks: an id that reappears after a boundary would silently absorb a
still-running rollout's labels and inflate coverage with no visible symptom. Revealing
labels per prefix instead would use information a deployed guard does not have, so we
report only the rollout-safe protocol.

\subsection{Base predictor and features}
\label{app:predictor}

The point and quantile predictors are LightGBM models \cite{ke2017lightgbm} with
\texttt{n\_estimators}${=}200$, \texttt{learning\_rate}${=}0.05$,
\texttt{num\_leaves}${=}15$ and \texttt{min\_child\_samples}${=}10$. CQR fits two
quantile regressors, at $\alpha/2$ and $1{-}\alpha/2$; the absolute-residual baseline
and the interval-driven stopping variant use a squared-error regressor; the feasibility
gate is an LGBM classifier with the same settings. No hyperparameter search was run:
these are the defaults fixed at the start of the project and used unchanged for every
experiment, baseline included.

The feature matrix (\texttt{features.build\_xy}) is: the turn index $k$; cumulative
spend; burn rate ($\mathrm{spend}/k$); the last, mean, standard deviation, minimum and
maximum per-turn spend; the mean of the last three turns; and the difference between
the last turn and the mean. When the agent emits its own interval, four columns are
added (the two endpoints, the midpoint and the width), and a missing estimate is left
as \texttt{NaN}, which LightGBM handles natively. Quantities known only at rollout end
(final totals, success, feasibility) are listed in a leak set and asserted absent from
the matrix at build time, so the exclusion is enforced rather than intended.

\subsection{Conformal methods and their settings}
\label{app:conformal}

Every method uses the finite-sample corrected quantile: with $n$ calibration scores the
level is $\lceil (n+1)(1{-}\alpha) \rceil / n$, read with higher interpolation,
returning $+\infty$ when that level exceeds 1. All intervals are clipped below at zero,
since a remaining budget cannot be negative. $\alpha{=}0.10$ throughout.

split-CQR uses the score $\max(\hat f_{\mathrm{lo}}-y,\ y-\hat f_{\mathrm{hi}})$ and
one global $\hat q$; it is the offline default. Mondrian-CQR applies the same score
with one $\hat q$ per turn-depth bucket, split at $k{=}1$, $k{=}2$ and $k{\geq}3$; a
bucket holding fewer than 10 calibration scores falls back to the global $\hat q$.

ACI updates the working level by $\alpha_{t+1} = \alpha_t + \gamma(\alpha -
\mathrm{err}_t)$. $\gamma$ is selected on calibration data only: the calibration
rollouts are split in half, ACI is streamed over the second half under the same
rollout-safe revelation used at test time, and the $\gamma \in \{0.005, 0.01, 0.02,
0.05\}$ whose coverage lands closest to $1{-}\alpha$ is kept. One deliberate deviation
from canonical ACI: where the canonical algorithm emits the empty set once $\alpha_t
\geq 1$, we clamp the quantile lookup at 0.999 and emit a finite interval. The
deviation is conservative in both coverage and the error feedback, and we state it
rather than diverge silently.

windowed-CQR keeps a rolling window of the 100 most recent conformity scores,
warm-started from calibration, and reads the interval half-width off that window's
empirical quantile, so the score distribution itself tracks drift rather than
reweighting a frozen calibration set. SPCI \cite{xu2023spci} is the inspiration, not
the implemented method.

DtACI aggregates ACI experts at $\gamma \in \{0.005, 0.01, 0.02, 0.05, 0.1, 0.2\}$
with learning rate $\eta = 2.0$ and uniform mixing $\sigma = 0.01$. After each revealed
label, experts are scored by pinball loss against $\beta_t$, the largest miscoverage
level whose quantile would still have covered the observed score, and the working level
is the weight-averaged $\alpha_t$.

\subsection{Baselines}
\label{app:baselines}

Seven comparators run on the same splits and seeds as the guard.

\begin{enumerate}
  \item \textbf{Agent self-report.} BAGEN's own per-sample indicators of whether the
  agent's interval contained the truth, read from the official estimation records. No
  calibration is applied.
  \item \textbf{Constant-quantile.} Feature-free: the empirical $[\alpha/2,
  1{-}\alpha/2]$ interval of the calibration labels, identical at every test prefix.
  \item \textbf{Split-absres.} A LightGBM point prediction plus or minus the calibrated
  absolute-residual quantile.
  \item \textbf{Conformalized-agent.} The agent's own interval midpoint used as the
  point predictor and conformalized on $|y - \mathrm{mid}|$.
  \item \textbf{Guard+self-est.} split-CQR on the guard features augmented with the
  agent's estimate columns.
  \item \textbf{Frozen split-CQR.} split-CQR calibrated on the source environment or
  model and applied unchanged to the target stream.
  \item \textbf{Hard caps.} Abort once spend so far exceeds a fixed cap. Dollar caps
  are swept at the 20th, 35th, 50th, 65th and 80th percentiles of realised run totals,
  and week caps at $0.5\times$ the 25th, $0.75\times$ the 50th and $1.0\times$ the 75th
  percentile of the task's time budget.
\end{enumerate}

Every baseline sees the same training split, the same calibration split, the same test
stream and the same seeds as the guard. BAGEN's SFT+RL estimator is not included as a
baseline: it is a fine-tuned agent rather than an external wrapper, and re-running it
would require training the underlying LLM, which is exactly what this work avoids.

\subsection{Metrics}
\label{app:metrics}

Coverage is the fraction of test prefixes whose interval contains the true remaining
requirement. Normalised width is the interval width divided by the central 80\% range
of the labels, $q_{0.9}{-}q_{0.1}$, which makes environments whose token scales differ
by orders of magnitude comparable; the median width is used wherever a method can emit
unbounded intervals, and the across-seed summary is reported as ``--'' where at least
one seed's median is itself infinite. Unbounded fraction is the share of predictions
with an infinite upper endpoint, and it appears beside every coverage number.

Confidence intervals come from a cluster bootstrap that resamples whole rollouts rather
than prefixes, with $B = 500$ resamples; per-sample binomial intervals are too tight
here, because the prefixes of one rollout are correlated. Across seeds, the per-seed
bootstrap draws are pooled and the 5th and 95th percentiles of the pooled distribution
are reported as a 90\% interval; these are recorded in the released result tables
rather than plotted. The $\pm$ shown in Table~\ref{tab:a2} and the error bars in
Figure~\ref{fig:coverage} are both the standard deviation across the five seeds.

For the stopping rule, the false-abort rate is the fraction of feasible test rollouts
that were stopped, and the doomed-savings fraction is the spend avoided on rollouts
that could not finish divided by the spend those rollouts would otherwise have
incurred. Cost-of-Pass is total spend under a policy divided by the number of runs that
still pass \cite{erol2026costofpass}. All economic results in this paper are Warehouse,
which is priced natively in simulated dollars; token-environment economics appear only
in Table~\ref{tab:abort}.

\begin{table}[t]
  \caption{Full within-run results: every conformal method lands within 0.05 of the
  0.90 target except the feature-free constant-quantile baseline on Search-R1, where
  the agent's own intervals miss by 0.39--0.67; windowed-CQR gives the narrowest
  bounded intervals on Sokoban and the second-narrowest on SWE-bench, while
  constant-quantile is sharpest on Search-R1. Claude Opus 4.7, nominal 0.90, mean
  $\pm$ standard deviation over five rollout-level split seeds.
  Normalised width is interval width divided by the central 80\% range of the labels
  (lower is sharper). ``Unb.'' is the percentage of predictions with an infinite upper
  endpoint; width is shown as ``--'' where at least one seed produced a majority of
  unbounded intervals. $n_{\mathrm{test}}$ is test prefixes per seed.}
  \label{tab:a2}
  \centering
  \small
  \begin{tabular}{llcccc}
    \toprule
    Environment & Method & Coverage & Norm.\ width ($\downarrow$) & Unb.\ (\%) & $n_{\mathrm{test}}$ \\
    \midrule
    Sokoban & agent self-report & 0.51 & -- & -- & -- \\
     & constant-quantile & 0.88 $\pm$ 0.05 & 1.10 & 0 & 87 \\
     & split-absres & 0.89 $\pm$ 0.03 & 1.12 & 0 & 87 \\
     & split-CQR & 0.92 $\pm$ 0.03 & 0.88 & 0 & 87 \\
     & conformalized-agent & 0.93 $\pm$ 0.05 & 1.44 & 0 & 87 \\
     & guard+self-est & 0.91 $\pm$ 0.04 & 0.91 & 0 & 87 \\
     & ACI-CQR & 0.92 $\pm$ 0.02 & 0.88 & 2 & 87 \\
     & windowed-CQR & 0.92 $\pm$ 0.02 & 0.86 & 0 & 87 \\
     & DtACI & 0.92 $\pm$ 0.02 & 0.88 & 0 & 87 \\
    \midrule
    Search-R1 & agent self-report & 0.23 & -- & -- & -- \\
     & constant-quantile & 0.77 $\pm$ 0.17 & 0.90 & 0 & 24 \\
     & split-absres & 0.94 $\pm$ 0.06 & 1.60 & 0 & 24 \\
     & split-CQR & 0.90 $\pm$ 0.15 & 1.14 & 0 & 24 \\
     & conformalized-agent & 0.90 $\pm$ 0.10 & 1.96 & 0 & 24 \\
     & guard+self-est & 0.85 $\pm$ 0.17 & 1.07 & 0 & 24 \\
     & ACI-CQR & 0.90 $\pm$ 0.15 & 1.14 & 0 & 24 \\
     & windowed-CQR & 0.92 $\pm$ 0.10 & 1.10 & 0 & 24 \\
     & DtACI & 0.94 $\pm$ 0.06 & -- & 16 & 24 \\
    \midrule
    SWE-bench & agent self-report & 0.35 & -- & -- & -- \\
     & constant-quantile & 0.92 $\pm$ 0.03 & 2.87 & 0 & 161 \\
     & split-absres & 0.95 $\pm$ 0.09 & 2.56 & 0 & 161 \\
     & split-CQR & 0.96 $\pm$ 0.10 & 2.11 & 0 & 161 \\
     & conformalized-agent & 0.97 $\pm$ 0.07 & 2.23 & 0 & 161 \\
     & guard+self-est & 0.96 $\pm$ 0.09 & 1.97 & 0 & 161 \\
     & ACI-CQR & 0.94 $\pm$ 0.09 & -- & 14 & 161 \\
     & windowed-CQR & 0.87 $\pm$ 0.08 & 1.29 & 0 & 161 \\
     & DtACI & 0.87 $\pm$ 0.07 & 1.22 & 0 & 161 \\
    \bottomrule
  \end{tabular}
\end{table}

\subsection{Multi-resource and censored calibration}
\label{app:joint}

Warehouse binds time in weeks, capacity in item-weeks and cost in simulated dollars at
once. Independent calibration sets each margin at $\alpha$; Bonferroni sets each at
$\alpha/3$, which restores a joint guarantee of $1{-}\alpha$ by the union bound and is
the primary multi-resource calibration scheme. The empirical-copula variant shares one
rank across resources and is tighter, but that rank is selected on the same calibration
scores it then thresholds, so it carries no guarantee and is reported as an observation
only. Under a model swap, an adaptive variant recomputes both the margins and the
shared rank from a rolling window of recent three-resource error vectors, updated only
once a rollout has finished.

For censored runs the one-sided upper bound is calibrated five ways
(\texttt{censored.upper\_qhat}): drop; infinity; bound; ipcw; and resource\_bound. A
Kaplan--Meier estimator \cite{kaplan1958nonparametric} is a diagnostic rather than a
sixth option: it returns $+\infty$ exactly when the upper tail is entirely censored.
Table~\ref{tab:b4} reports the five calibrations together with the Kaplan--Meier
diagnostic. Censored-stratum coverage for capacity and cost is coverage against an
approximate interpolated label (12--27\% median relative error), not against ground
truth.

The feasibility gate is an LGBM classifier on the same prefix features. It issues
intervals only on predicted-feasible prefixes, so the resulting guarantee is selective.
The gate is swept at its own prefix-level tolerance $\delta_{\mathrm{gate}} \in
\{0.02, 0.05, 0.10, 0.20, 0.30\}$, which is a different quantity from the rollout-level
false-abort tolerance $\delta$ of \S\ref{sec:econ}; Table~\ref{tab:b3} reports the full
sweep.

\subsection{Stopping rule and economics}
\label{app:stopping}

The reported rule aborts at the first turn whose feasibility score falls below $\tau$.
Because the rule aborts a rollout exactly when that rollout's minimum prefix score
crosses $\tau$, $\tau$ is calibrated at the rollout level: it is set to the score at
rank $\lfloor (n+1)\delta \rfloor$ among the sorted minimum scores of the $n$ feasible
calibration rollouts, which bounds by $\delta$ the probability of aborting a run that
would have finished. $\delta$ is swept over $\{0.02, 0.05, 0.10, 0.15, 0.20, 0.30,
0.40\}$ for the frontier; $\delta{=}0.10$ is the operating point quoted in
\S\ref{sec:econ} and the abstract.

An interval-driven variant, which aborts when spend so far plus the calibrated lower
bound exceeds the budget, is implemented in the released code and calibrated the same
way at rollout level. It is not the rule behind the reported Warehouse economics, and
we do not claim the two are equivalent. The economic numbers are counterfactual
replays, not live measurements: each completed run is truncated where the guard would
have aborted it and the avoided spend is summed.

\subsection{Data, environment, and reproduction}
\label{app:repro}

All experiments use logged rollouts from BAGEN \cite{lin2026bagen}, fetched from its
public Hugging Face release at the pinned revision
\texttt{58189576e54b675fdd0e1d6c1c9f189c2992732f}. Its dataset card states no licence,
so we use the data as publicly released for research and credit the originators rather
than asserting terms the release does not state. The extractor reproduces BAGEN's
turn-excluding-history token accounting exactly, with 0 of 355 mismatches against the
official per-sample ground truth on the reference run.

The whole pipeline runs inside a container whose base image is pinned by digest
(\texttt{python:3.12-slim-bookworm}), with numpy 2.5.1, pandas 3.0.3, LightGBM 4.6.0,
scikit-learn 1.9.0, scipy 1.18.0 and matplotlib 3.11.0 resolved from a hash-locked
requirements file. LightGBM is MIT-licensed; numpy, pandas, scipy and scikit-learn are
BSD-3-Clause; matplotlib uses a PSF-style licence. The guard is CPU-only: no GPU, no
model weights, and no access to logits or environment internals. No LLM is trained,
fine-tuned or queried by us at any stage. A single command regenerates every figure and
table in the repository and then machine-verifies the paper against what it just
produced, exiting non-zero if any reported claim no longer matches.

\section*{NeurIPS Paper Checklist}

\begin{enumerate}

\item {\bf Claims}
    \item[] Question: Do the main claims made in the abstract and introduction accurately reflect the paper's contributions and scope?
    \item[] Answer: \answerYes{}
    \item[] Justification: Every headline number in the abstract and Section~\ref{sec:intro} is produced by an experiment in this paper. The abstract distinguishes interval calibration from abort economics, states that the 86.3\% savings and Cost-of-Pass figures are Warehouse-only and do not use remaining-budget intervals, and reports token-environment abort results at the same $\delta$. Scope limits are stated throughout: Claude Opus 4.7 for within-run coverage, 0.95 joint coverage in-distribution only (0.71 under model swap), and counterfactual replay rather than live deployment.
    \item[] Guidelines:
    \begin{itemize}
        \item The answer \answerNA{} means that the abstract and introduction do not include the claims made in the paper.
        \item The abstract and/or introduction should clearly state the claims made, including the contributions made in the paper and important assumptions and limitations. A \answerNo{} or \answerNA{} answer to this question will not be perceived well by the reviewers.
        \item The claims made should match theoretical and experimental results, and reflect how much the results can be expected to generalize to other settings.
        \item It is fine to include aspirational goals as motivation as long as it is clear that these goals are not attained by the paper.
    \end{itemize}

\item {\bf Limitations}
    \item[] Question: Does the paper discuss the limitations of the work performed by the authors?
    \item[] Answer: \answerYes{}
    \item[] Justification: Section~\ref{sec:limits} states that the joint multi-resource result rests entirely on one environment and does not survive a model swap, that the breadth counts come from a single split rather than a seed average, that marginal coverage can hide under-coverage in slices such as early turns, and that the economic figures are counterfactual replays rather than live measurements. Appendix~\ref{app:protocol} adds the assumption-level caveats and Table~\ref{tab:a1} marks which components carry no finite-sample guarantee.
    \item[] Guidelines:
    \begin{itemize}
        \item The answer \answerNA{} means that the paper has no limitation while the answer \answerNo{} means that the paper has limitations, but those are not discussed in the paper.
        \item The authors are encouraged to create a separate ``Limitations'' section in their paper.
        \item The paper should point out any strong assumptions and how robust the results are to violations of these assumptions (e.g., independence assumptions, noiseless settings, model well-specification, asymptotic approximations only holding locally). The authors should reflect on how these assumptions might be violated in practice and what the implications would be.
        \item The authors should reflect on the scope of the claims made, e.g., if the approach was only tested on a few datasets or with a few runs. In general, empirical results often depend on implicit assumptions, which should be articulated.
        \item The authors should reflect on the factors that influence the performance of the approach. For example, a facial recognition algorithm may perform poorly when image resolution is low or images are taken in low lighting. Or a speech-to-text system might not be used reliably to provide closed captions for online lectures because it fails to handle technical jargon.
        \item The authors should discuss the computational efficiency of the proposed algorithms and how they scale with dataset size.
        \item If applicable, the authors should discuss possible limitations of their approach to address problems of privacy and fairness.
        \item While the authors might fear that complete honesty about limitations might be used by reviewers as grounds for rejection, a worse outcome might be that reviewers discover limitations that aren't acknowledged in the paper. The authors should use their best judgment and recognize that individual actions in favor of transparency play an important role in developing norms that preserve the integrity of the community. Reviewers will be specifically instructed to not penalize honesty concerning limitations.
    \end{itemize}

\item {\bf Theory assumptions and proofs}
    \item[] Question: For each theoretical result, does the paper provide the full set of assumptions and a complete (and correct) proof?
    \item[] Answer: \answerNA{}
    \item[] Justification: The paper states and proves no new theoretical results. Every guarantee is imported from prior work and Table~\ref{tab:a1} records what is guaranteed for each component.
    \item[] Guidelines:
    \begin{itemize}
        \item The answer \answerNA{} means that the paper does not include theoretical results.
        \item All the theorems, formulas, and proofs in the paper should be numbered and cross-referenced.
        \item All assumptions should be clearly stated or referenced in the statement of any theorems.
        \item The proofs can either appear in the main paper or the supplemental material, but if they appear in the supplemental material, the authors are encouraged to provide a short proof sketch to provide intuition.
        \item Inversely, any informal proof provided in the core of the paper should be complemented by formal proofs provided in appendix or supplemental material.
        \item Theorems and Lemmas that the proof relies upon should be properly referenced.
    \end{itemize}

\item {\bf Experimental result reproducibility}
    \item[] Question: Does the paper fully disclose all the information needed to reproduce the main experimental results of the paper to the extent that it affects the main claims and/or conclusions of the paper (regardless of whether the code and data are provided or not)?
    \item[] Answer: \answerYes{}
    \item[] Justification: Appendix~\ref{app:protocol} specifies the split procedure, rollout-safe protocol, LightGBM settings, feature list, every conformal variant, all seven baselines, every metric, and the pinned data revision and software environment.
    \item[] Guidelines:
    \begin{itemize}
        \item The answer \answerNA{} means that the paper does not include experiments.
        \item If the paper includes experiments, a \answerNo{} answer to this question will not be perceived well by the reviewers: Making the paper reproducible is important, regardless of whether the code and data are provided or not.
        \item If the contribution is a dataset and\slash or model, the authors should describe the steps taken to make their results reproducible or verifiable.
        \item Depending on the contribution, reproducibility can be accomplished in various ways. For example, if the contribution is a novel architecture, describing the architecture fully might suffice, or if the contribution is a specific model and empirical evaluation, it may be necessary to either make it possible for others to replicate the model with the same dataset, or provide access to the model. In general. releasing code and data is often one good way to accomplish this, but reproducibility can also be provided via detailed instructions for how to replicate the results, access to a hosted model (e.g., in the case of a large language model), releasing of a model checkpoint, or other means that are appropriate to the research performed.
        \item While NeurIPS does not require releasing code, the conference does require all submissions to provide some reasonable avenue for reproducibility, which may depend on the nature of the contribution. For example
        \begin{enumerate}
            \item If the contribution is primarily a new algorithm, the paper should make it clear how to reproduce that algorithm.
            \item If the contribution is primarily a new model architecture, the paper should describe the architecture clearly and fully.
            \item If the contribution is a new model (e.g., a large language model), then there should either be a way to access this model for reproducing the results or a way to reproduce the model (e.g., with an open-source dataset or instructions for how to construct the dataset).
            \item We recognize that reproducibility may be tricky in some cases, in which case authors are welcome to describe the particular way they provide for reproducibility. In the case of closed-source models, it may be that access to the model is limited in some way (e.g., to registered users), but it should be possible for other researchers to have some path to reproducing or verifying the results.
        \end{enumerate}
    \end{itemize}

\item {\bf Open access to data and code}
    \item[] Question: Does the paper provide open access to the data and code, with sufficient instructions to faithfully reproduce the main experimental results, as described in supplemental material?
    \item[] Answer: \answerNo{}
    \item[] Justification: The data is fully open (public BAGEN rollouts at a pinned revision). The code is not attached to this submission; the full pipeline will be released under an MIT licence on acceptance, together with a digest-pinned Dockerfile and a single entry point that regenerates and verifies every reported number. Appendix~\ref{app:protocol} specifies everything needed for an independent reimplementation in the meantime.
    \item[] Guidelines:
    \begin{itemize}
        \item The answer \answerNA{} means that paper does not include experiments requiring code.
        \item Please see the NeurIPS code and data submission guidelines (\url{https://neurips.cc/public/guides/CodeSubmissionPolicy}) for more details.
        \item While we encourage the release of code and data, we understand that this might not be possible, so \answerNo{} is an acceptable answer. Papers cannot be rejected simply for not including code, unless this is central to the contribution (e.g., for a new open-source benchmark).
        \item The instructions should contain the exact command and environment needed to run to reproduce the results. See the NeurIPS code and data submission guidelines (\url{https://neurips.cc/public/guides/CodeSubmissionPolicy}) for more details.
        \item The authors should provide instructions on data access and preparation, including how to access the raw data, preprocessed data, intermediate data, and generated data, etc.
        \item The authors should provide scripts to reproduce all experimental results for the new proposed method and baselines. If only a subset of experiments are reproducible, they should state which ones are omitted from the script and why.
        \item At submission time, to preserve anonymity, the authors should release anonymized versions (if applicable).
        \item Providing as much information as possible in supplemental material (appended to the paper) is recommended, but including URLs to data and code is permitted.
    \end{itemize}

\item {\bf Experimental setting/details}
    \item[] Question: Does the paper specify all the training and test details (e.g., data splits, hyperparameters, how they were chosen, type of optimizer) necessary to understand the results?
    \item[] Answer: \answerYes{}
    \item[] Justification: Section~\ref{sec:setup} and Table~\ref{tab:envs} give the environments and budgets; Appendix~\ref{app:protocol} gives the rest, including the exact sweep grids for hard-cap baselines, the feasibility gate, and $\delta$.
    \item[] Guidelines:
    \begin{itemize}
        \item The answer \answerNA{} means that the paper does not include experiments.
        \item The experimental setting should be presented in the core of the paper to a level of detail that is necessary to appreciate the results and make sense of them.
        \item The full details can be provided either with the code, in appendix, or as supplemental material.
    \end{itemize}

\item {\bf Experiment statistical significance}
    \item[] Question: Does the paper report error bars suitably and correctly defined or other appropriate information about the statistical significance of the experiments?
    \item[] Answer: \answerYes{}
    \item[] Justification: Every headline number is a mean over five rollout-level split seeds. Figure~\ref{fig:coverage} carries across-seed standard-deviation error bars; Table~\ref{tab:a2} reports the same quantity as a $\pm$ beside each coverage value. Appendix~\ref{app:metrics} specifies the cluster bootstrap ($B = 500$, resampling whole rollouts).
    \item[] Guidelines:
    \begin{itemize}
        \item The answer \answerNA{} means that the paper does not include experiments.
        \item The authors should answer \answerYes{} if the results are accompanied by error bars, confidence intervals, or statistical significance tests, at least for the experiments that support the main claims of the paper.
        \item The factors of variability that the error bars are capturing should be clearly stated (for example, train/test split, initialization, random drawing of some parameter, or overall run with given experimental conditions).
        \item The method for calculating the error bars should be explained (closed form formula, call to a library function, bootstrap, etc.)
        \item The assumptions made should be given (e.g., Normally distributed errors).
        \item It should be clear whether the error bar is the standard deviation or the standard error of the mean.
        \item It is OK to report 1-sigma error bars, but one should state it. The authors should preferably report a 2-sigma error bar than state that they have a 96\% CI, if the hypothesis of Normality of errors is not verified.
        \item For asymmetric distributions, the authors should be careful not to show in tables or figures symmetric error bars that would yield results that are out of range (e.g., negative error rates).
        \item If error bars are reported in tables or plots, the authors should explain in the text how they were calculated and reference the corresponding figures or tables in the text.
    \end{itemize}

\item {\bf Experiments compute resources}
    \item[] Question: For each experiment, does the paper provide sufficient information on the computer resources (type of compute workers, memory, time of execution) needed to reproduce the experiments?
    \item[] Answer: \answerYes{}
    \item[] Justification: Appendix~\ref{app:repro} states that the entire pipeline is CPU-only: no GPU, no model weights, and no LLM training or querying. All experiments fit gradient-boosted trees on tables of at most a few thousand rows inside one digest-pinned container.
    \item[] Guidelines:
    \begin{itemize}
        \item The answer \answerNA{} means that the paper does not include experiments.
        \item The paper should indicate the type of compute workers CPU or GPU, internal cluster, or cloud provider, including relevant memory and storage.
        \item The paper should provide the amount of compute required for each of the individual experimental runs as well as estimate the total compute.
        \item The paper should disclose whether the full research project required more compute than the experiments reported in the paper (e.g., preliminary or failed experiments that didn't make it into the paper).
    \end{itemize}

\item {\bf Code of ethics}
    \item[] Question: Does the research conducted in the paper conform, in every respect, with the NeurIPS Code of Ethics \url{https://neurips.cc/public/EthicsGuidelines}?
    \item[] Answer: \answerYes{}
    \item[] Justification: We have reviewed the NeurIPS Code of Ethics and believe this work conforms. The research uses only an existing, publicly released dataset of logged agent rollouts; it involves no human subjects and no new model release.
    \item[] Guidelines:
    \begin{itemize}
        \item The answer \answerNA{} means that the authors have not reviewed the NeurIPS Code of Ethics.
        \item If the authors answer \answerNo, they should explain the special circumstances that require a deviation from the Code of Ethics.
        \item The authors should make sure to preserve anonymity (e.g., if there is a special consideration due to laws or regulations in their jurisdiction).
    \end{itemize}

\item {\bf Broader impacts}
    \item[] Question: Does the paper discuss both potential positive societal impacts and negative societal impacts of the work performed?
    \item[] Answer: \answerYes{}
    \item[] Justification: Section~\ref{sec:limits} contains an explicit broader-impact paragraph on wasted compute versus false-abort risk, with $\delta$ calibration as the mitigation.
    \item[] Guidelines:
    \begin{itemize}
        \item The answer \answerNA{} means that there is no societal impact of the work performed.
        \item If the authors answer \answerNA{} or \answerNo, they should explain why their work has no societal impact or why the paper does not address societal impact.
        \item Examples of negative societal impacts include potential malicious or unintended uses (e.g., disinformation, generating fake profiles, surveillance), fairness considerations (e.g., deployment of technologies that could make decisions that unfairly impact specific groups), privacy considerations, and security considerations.
        \item The conference expects that many papers will be foundational research and not tied to particular applications, let alone deployments. However, if there is a direct path to any negative applications, the authors should point it out. For example, it is legitimate to point out that an improvement in the quality of generative models could be used to generate Deepfakes for disinformation. On the other hand, it is not needed to point out that a generic algorithm for optimizing neural networks could enable people to train models that generate Deepfakes faster.
        \item The authors should consider possible harms that could arise when the technology is being used as intended and functioning correctly, harms that could arise when the technology is being used as intended but gives incorrect results, and harms following from (intentional or unintentional) misuse of the technology.
        \item If there are negative societal impacts, the authors could also discuss possible mitigation strategies (e.g., gated release of models, providing defenses in addition to attacks, mechanisms for monitoring misuse, mechanisms to monitor how a system learns from feedback over time, improving the efficiency and accessibility of ML).
    \end{itemize}

\item {\bf Safeguards}
    \item[] Question: Does the paper describe safeguards that have been put in place for responsible release of data or models that have a high risk for misuse (e.g., pre-trained language models, image generators, or scraped datasets)?
    \item[] Answer: \answerNA{}
    \item[] Justification: We release no model weights, no pretrained language model, and no scraped data. The artifact is CPU-only calibration code whose only effect is to stop an agent's own run earlier.
    \item[] Guidelines:
    \begin{itemize}
        \item The answer \answerNA{} means that the paper poses no such risks.
        \item Released models that have a high risk for misuse or dual-use should be released with necessary safeguards to allow for controlled use of the model, for example by requiring that users adhere to usage guidelines or restrictions to access the model or implementing safety filters.
        \item Datasets that have been scraped from the Internet could pose safety risks. The authors should describe how they avoided releasing unsafe images.
        \item We recognize that providing effective safeguards is challenging, and many papers do not require this, but we encourage authors to take this into account and make a best faith effort.
    \end{itemize}

\item {\bf Licenses for existing assets}
    \item[] Question: Are the creators or original owners of assets (e.g., code, data, models), used in the paper, properly credited and are the license and terms of use explicitly mentioned and properly respected?
    \item[] Answer: \answerYes{}
    \item[] Justification: BAGEN~\citep{lin2026bagen} is cited and pinned to an exact revision; its dataset card states no licence, so we assert none on its behalf and credit the originators. Appendix~\ref{app:repro} lists the software stack with exact versions and licences (LightGBM MIT; numpy, pandas, scipy, scikit-learn BSD-3-Clause; matplotlib PSF-style).
    \item[] Guidelines:
    \begin{itemize}
        \item The answer \answerNA{} means that the paper does not use existing assets.
        \item The authors should cite the original paper that produced the code package or dataset.
        \item The authors should state which version of the asset is used and, if possible, include a URL.
        \item The name of the license (e.g., CC-BY 4.0) should be included for each asset.
        \item For scraped data from a particular source (e.g., website), the copyright and terms of service of that source should be provided.
        \item If assets are released, the license, copyright information, and terms of use in the package should be provided. For popular datasets, \url{paperswithcode.com/datasets} has curated licenses for some datasets. Their licensing guide can help determine the license of a dataset.
        \item For existing datasets that are re-packaged, both the original license and the license of the derived asset (if it has changed) should be provided.
        \item If this information is not available online, the authors are encouraged to reach out to the asset's creators.
    \end{itemize}

\item {\bf New assets}
    \item[] Question: Are new assets introduced in the paper well documented and is the documentation provided alongside the assets?
    \item[] Answer: \answerNA{}
    \item[] Justification: No new asset is released with this submission. The implementation will be released on acceptance; no new dataset is created.
    \item[] Guidelines:
    \begin{itemize}
        \item The answer \answerNA{} means that the paper does not release new assets.
        \item Researchers should communicate the details of the dataset\slash code\slash model as part of their submissions via structured templates. This includes details about training, license, limitations, etc.
        \item The paper should discuss whether and how consent was obtained from people whose asset is used.
        \item At submission time, remember to anonymize your assets (if applicable). You can either create an anonymized URL or include an anonymized zip file.
    \end{itemize}

\item {\bf Crowdsourcing and research with human subjects}
    \item[] Question: For crowdsourcing experiments and research with human subjects, does the paper include the full text of instructions given to participants and screenshots, if applicable, as well as details about compensation (if any)?
    \item[] Answer: \answerNA{}
    \item[] Justification: All experiments are offline analyses of pre-existing logged agent rollouts.
    \item[] Guidelines:
    \begin{itemize}
        \item The answer \answerNA{} means that the paper does not involve crowdsourcing nor research with human subjects.
        \item Including this information in the supplemental material is fine, but if the main contribution of the paper involves human subjects, then as much detail as possible should be included in the main paper.
        \item According to the NeurIPS Code of Ethics, workers involved in data collection, curation, or other labor should be paid at least the minimum wage in the country of the data collector.
    \end{itemize}

\item {\bf Institutional review board (IRB) approvals or equivalent for research with human subjects}
    \item[] Question: Does the paper describe potential risks incurred by study participants, whether such risks were disclosed to the subjects, and whether Institutional Review Board (IRB) approvals (or an equivalent approval/review based on the requirements of your country or institution) were obtained?
    \item[] Answer: \answerNA{}
    \item[] Justification: No human subjects were involved.
    \item[] Guidelines:
    \begin{itemize}
        \item The answer \answerNA{} means that the paper does not involve crowdsourcing nor research with human subjects.
        \item Depending on the country in which research is conducted, IRB approval (or equivalent) may be required for any human subjects research. If you obtained IRB approval, you should clearly state this in the paper.
        \item We recognize that the procedures for this may vary significantly between institutions and locations, and we expect authors to adhere to the NeurIPS Code of Ethics and the guidelines for their institution.
        \item For initial submissions, do not include any information that would break anonymity (if applicable), such as the institution conducting the review.
    \end{itemize}

\item {\bf Declaration of LLM usage}
    \item[] Question: Does the paper describe the usage of LLMs if it is an important, original, or non-standard component of the core methods in this research? Note that if the LLM is used only for writing, editing, or formatting purposes and does \emph{not} impact the core methodology, scientific rigor, or originality of the research, declaration is not required.
    \item[] Answer: \answerNA{}
    \item[] Justification: No LLM is a component of the core method. We neither train, fine-tune nor query any LLM in this work; the seven agents named in Section~\ref{sec:setup} are the sources of the pre-existing logged rollouts we analyse, not components of the method.
    \item[] Guidelines:
    \begin{itemize}
        \item The answer \answerNA{} means that the core method development in this research does not involve LLMs as any important, original, or non-standard components.
        \item Please refer to our LLM policy in the NeurIPS handbook for what should or should not be described.
    \end{itemize}

\end{enumerate}

\end{document}
